\section*{Chapitre \ref{ch:s2:hedging-a-structured-products-book} --- Couverture d'un livre de produits structurés}

\begin{solution}{exo:s2:hedging-a-structured-products-book:1}
$(100 - 99.12)/100 \times \text{\euro}400$ millions $\approx$ 3.52~millions d'euros.
\end{solution}

\begin{solution}{exo:s2:hedging-a-structured-products-book:2}
$9.66 / 5 = 1.93$ points~; la banque en verse en moyenne la moitié sur le vega qu'elle vend~: $9.66 \times 1.93 / 2 \approx$ 9.33~millions d'euros.
\end{solution}

\begin{solution}{exo:s2:hedging-a-structured-products-book:3}
La banque achète 228~millions d'euros de l'indice à mesure qu'il recule de 100~\% à 55~\%~: la couverture passe de 106~millions d'euros à 334~millions d'euros.
\end{solution}

\begin{solution}{exo:s2:hedging-a-structured-products-book:4}
L'investisseur est vendeur d'un put qui se déclenche sous le niveau de protection, et l'autocall plafonne la hausse qu'il pourrait tirer des coupons~; l'émetteur détient
l'autre côté. Un put prend de la valeur quand la volatilité monte, donc l'engagement de l'émetteur baisse~: il est long en volatilité.
\end{solution}

\begin{solution}{exo:s2:hedging-a-structured-products-book:5}
Le worst-of verse selon le moins bon de deux indices. Plus les deux indices évoluent de façon semblable, moins l'un d'eux a de chances de finir loin sous l'autre, si bien que la
note vaut davantage pour l'investisseur~; l'émetteur, vendeur de la note, perd lorsque la corrélation augmente.
\end{solution}

\begin{solution}{exo:s2:hedging-a-structured-products-book:6}
Près des barrières, la taille de la couverture varie le plus vite avec l'indice. Sous le niveau où elle culmine, l'émetteur doit vendre l'indice à mesure qu'il baisse, ce qui accentue
la baisse, et lorsque de nombreuses notes partagent les mêmes barrières et les mêmes sous-jacents, comme l'a décrit la BIS pour la Corée, les ventes sont importantes par rapport au marché.
\end{solution}

\begin{solution}{exo:s2:hedging-a-structured-products-book:7}
Entre 50~\% et 55~\% du strike, soit environ 52~\% par interpolation (vega $-0.47$ million à 50~\%, $+0.76$ million à 55~\%). Sous ce niveau, la protection est
probablement perdue, et la note est surtout une détention de l'indice~; une volatilité plus élevée augmente alors la probabilité que l'indice remonte au-dessus de 60~\% à l'échéance,
qui profite à l'investisseur et coûte à l'émetteur.
\end{solution}

\begin{solution}{exo:s2:hedging-a-structured-products-book:8}
La marge est de 8.31~millions d'euros pour 1~milliard d'euros~; recycler à lui seul le vega du livre coûte 4.67~millions d'euros à une profondeur de 10~millions d'euros par point, et
9.33~millions d'euros à 5~millions d'euros. Les expositions non couvrables et la concession nécessaire pour les transférer sont l'économie même du métier, non un
détail.
\end{solution}

\begin{solution}{pb:s2:hedging-a-structured-products-book:1}
\begin{enumerate}
\item Une sensibilité d'un livre de produits structurés à un paramètre ne se couvre qu'en partie avec des instruments vanille, mesurée en réévaluant le livre après choc du paramètre.
\item 400~millions d'euros sur A, 300~millions d'euros sur B, 300~millions d'euros en worst-of~; cinq ans, autocall annuel à 100~\%, coupons Phoenix de 7~\%
et 9~\% à 70~\%, protection de 60~\%~; volatilité de 20~\% avec skew, dividendes de 3~\% et corrélation de 0.7.
\item 8.31~millions d'euros~: les notes valent 99.12 et 99.29 pour 100.
\item Les autocallables sont largement vendus aux investisseurs particuliers coréens, couvrir ces produits peut exiger des trades déstabilisants près des barrières, et les émissions sur Samsung
et SK Hynix ont fortement augmenté au premier semestre 2026.
\item Vega 3.22, 2.42, 4.02, total 9.66~; skew 1.47, 1.10, 1.86, total 4.42~; dividendes 0.51, 0.38, 0.55, total 1.43~; corrélation $-1.05$ (\euro{}
millions).
\item Elle est longue du put que l'investisseur a vendu, et un put gagne à une volatilité plus élevée, à un skew plus pentu à la baisse et à des dividendes plus élevés.
\item L'investisseur en worst-of gagne lorsque les indices évoluent de concert~; l'émetteur détient l'autre côté.
\item De 3.22~millions d'euros au strike à un maximum de 4.45~millions d'euros près de 80~\%, devenant négatif sous environ 52~\%.
\item Transférer les expositions que le livre ne peut couvrir avec des instruments vanille~: vendre la volatilité et les dividendes de long terme, acheter de la corrélation.
\item Le vega via des options de long terme et des swaps de variance vendus aux hedge funds~; les dividendes via des contrats à terme sur dividendes et des swaps~; la corrélation en l'achetant
au moyen de trades de dispersion.
\item À des profondeurs de 5, 10 et 20~millions d'euros par point~: la volatilité implicite baisse de 1.93, 0.97 et 0.48 points~; la banque verse 9.33, 4.67 et 2.33~millions d'euros~;
un acheteur patient gagne 18.67, 9.33 et 4.67~millions d'euros.
\item Les acheteurs des expositions recyclées, payés par les émetteurs au moyen de la concession.
\item Position longue de 9.66~millions d'euros de vega, de 4.42~millions d'euros par 0.05 de skew et de 1.43~millions d'euros par 0.1 point de dividendes, position courte de 1.05~millions d'euros par 0.05 de
corrélation~; recycler le vega déplace la volatilité implicite de long terme de $-0.97$ point à une profondeur de 10~millions d'euros.
\item Sous 55~\%, la couverture se réduit à mesure que l'indice baisse~: la banque vend dans la baisse.
\item De même ordre de grandeur~: 4.67~millions d'euros, à comparer à une marge de 8.31~millions d'euros, à une profondeur de 10~millions d'euros.
\item Elles partagent la profondeur~: le mouvement et le coût augmentent avec l'offre totale.
\item La dispersion du chapitre~2, les \omterm{def:s2:skew-and-term-structure-relative-value:term}{trades de structure par terme} et de skew du chapitre~3, et les acheteurs de dividendes du chapitre~26.
\item Celui de l'acheteur de vega recyclé.
\item Sur la volatilité implicite de long terme cotée, avec ses fourchettes, en entrant après les vagues d'émissions et en conservant la position jusqu'au retour à la moyenne, y compris les années où il ne
revient pas.
\item Les expositions que lui laissent ses notes, recyclées auprès de quiconque les reprendra à un prix.
\end{enumerate}
\end{solution}

\begin{solution}{iq:s2:hedging-a-structured-products-book:1}
Position longue en volatilité (surtout de long terme), longue en skew à la baisse, longue en dividendes, courte en corrélation pour les notes worst-of, et risque de saut et de barrière qui croît à mesure que
le sous-jacent s'approche des niveaux de déclenchement.
\end{solution}

\begin{solution}{iq:s2:hedging-a-structured-products-book:2}
Réévaluer le livre après choc de la matrice de corrélation (uniformément et paire par paire), avec des nombres aléatoires communs, et exprimer le résultat sous forme de
vega de corrélation par paire et par échéance~; le confronter aux couvertures de dispersion du livre.
\end{solution}

\begin{solution}{iq:s2:hedging-a-structured-products-book:3}
Recalculer les expositions au nouveau niveau, puisque le vega et le delta y varient le plus vite~; réduire prudemment la couverture en delta au passage des barrières, acheter
une protection à la baisse si elle n'est pas trop chère, et coordonner avec les limites de risque avant que le marché n'aille plus loin~; s'attendre à ce que d'autres ayant les
mêmes livres fassent de même.
\end{solution}

\begin{solution}{iq:s2:hedging-a-structured-products-book:4}
De fortes baisses au passage des barrières, avec sursaut de la volatilité et de la corrélation, des baisses de dividendes, et un mouvement de saut à une date d'observation, chacun s'accompagnant
d'une liquidité réduite sur les marchés de couverture.
\end{solution}

\begin{solution}{iq:s2:hedging-a-structured-products-book:5}
Mutualiser les trajectoires simulées entre notes ayant les mêmes sous-jacents, utiliser les sensibilités adjointes (pathwise) plutôt que des chocs, réduire la grille temporelle aux
dates d'observation lorsque le payoff le permet, et exécuter le calcul sur de nombreux cœurs ou GPU.
\end{solution}

\begin{solution}{iq:s2:hedging-a-structured-products-book:6}
Vendre $x$ déplace le prix de $x/D$~; le prix marginal de la $v$-ième unité vaut $v/D$, si bien que le coût vaut $\int_0^V v/D\,dv = V^2/(2D)$. En $n$ tranches
avec un rétablissement complet du prix entre deux tranches, chacune coûte $(V/n)^2/(2D)$, et les $n$ coûtent $V^2/(2nD)$.
\end{solution}
